\documentclass[a4paper]{article}
\usepackage[margin=3cm]{geometry}

\usepackage{hyperref}
\hypersetup{
	colorlinks=true,
	linkcolor=blue,
	filecolor=magenta,
	urlcolor=cyan,
	hyperindex=true,
	linktocpage=true,%makes the page number be the link in the index
}

\title{Computational bio-engineering}

\author{Chaya2766}

\begin{document}
	\sffamily
	\sloppy
	
	\maketitle
	
	\vfill
	
	\begin{abstract}
		Approximation of the computational power required for in-software simulation of effects of biological engineering on reasonable timescales for the needs of the Stareater Expanse setting.
		
		\medskip
		
		2 levels of abstraction are inspected: cellular and molecular.
		
		\medskip
		
		Cellular-level abstraction is estimated to require 1.24e10 bits of memory and 6.075e19 bits/s of flip rate per cell for a realtime simulation. This is far too much for any useful simulation size to be performed within a person's cybernetic implants, but likely enough for larger computational clusters to make useful predictions on matters such as organ rejection.
		
		\medskip
		
		Molecular-level abstraction is estimated to require on average 3.6756e23 bits/s of flip rate for every atom in the simulation, corresponding to 1.106e38 bits/s for a typical fibroblast cell, and 3.687e46 bits/s for a whole milliliter of biological tissue. This limits the ability to perform such simulations with reasonable amounts of time, energy and hardware to biological systems below the size of single cells or at most small groups of cells.
	\end{abstract}
	
	%\noindent \rule{\linewidth}{1pt}
	\vfill
	
	\tableofcontents
	
	\pagebreak
	
	\section{Possible levels of abstraction}
	
	The required detail of simulation may differ, here are some main possibilities I foresee and how likely I rate them:
	
	\begin{enumerate}
		\item Can someone decide on the changes they want made to an organism and then elegantly solve for the genetic modifications needed to achieve them? My limited knowledge of how life works pushes me to believe that even the smallest change to an organism's genome will result in a cascade of differences showing up in interactions between that particular gene and every other existing gene, and in turn countless cascades of changes resulting from those slightly different interactions themselves interacting with one another to produce new results, making it a brute-force problem. Looking at it this way, the only simple solution that allows this may be a database that can tell you what genetic modification results in what changes to the organism in past experiments.
		
		\item Abstracting entire cells as the smallest unit needed to be simulated may be sufficient. I can't write this off as a possibility, especially in the case where one works with known, pre-tested gene combinations rather than exploring entirely new arrangements. It might be possible to produce an abstract model of a cell that behaves as expected within known boundaries, but I don't have a good idea of how one might approximate the computational power required for simulations running on this level.
		
		\item Simulating every molecule may be necessary for accurate results. This would not require any advanced understanding of how biological systems behave, but merely how the system is structured and how the individual molecules behave. To me this seems like the most promising approach in terms of asking "would it work?" as we already use molecule simulations to develop medicine in the modern day, the question is just what it would take in terms of computing.
		
		\item Subatomic details may be required for the system to function. Life has a tendency to incorporate every possible effect into it, looking only for whether or not it is useful rather than any other considerations. It is not out of question that effects such as quantum tunneling would have found themselves used by some system somewhere in biology, eg. removing the effect makes one enzyme somewhere work ever so slightly differently and through the messy web of relations that biological pathways are somehow breaks everything. In that case it would be necessary to simulate everything down to the level of subatomic particles or waves, which would be a huge step up on computational complexity.
	\end{enumerate}
	
	\medskip
	
	\noindent\rule{\linewidth}{1pt}
	
	\medskip
	
	For the needs of the Stareater Expanse setting, as of writing this document my stance is this:
	
	\begin{enumerate}
		\item Solving for necessary genetic changes or for the effects of said changes with simple calculations or lookup tables is possible and reliable for thoroughly tested gene combinations, but only works within that limited scope, and does not allow for predicting effects of new combinations ahead of time.
		
		\item Cell-level abstraction is available and can predict problems before implementing changes, but does not guarantee accuracy for new genetic combinations. This model would be useful for example in determining whether an artificial organ with slight genetic differences will be rejected by a patient's body or not, ie. medical environments working with already known genetics. It can however be used to search for promising pathways for novel genetic designs before they are proven with more computationally demanding simulations.
		
		\item molecular-level abstraction can reasonably guarantee accuracy for new genetic combinations. This model suffers in the ethics aspect, as it would be sufficient to fully simulate neurological activity in the case that embryos are simulated to validate reproductive stability of the new genome.
		
		\item subatomic-level simulation of cellular-scale biological systems to account for all physical effects is not plausible with classical computing at the kardashev-2 scale
	\end{enumerate}
	
	An additional consideration is that finding viable genetic modifications might involve running thousands to billions of simulations, but once found and recorded they could likely be verified with even just one simulation.
	
	\pagebreak
	
	\section{Cell-level abstraction}
	
	I suspect the complexity of the computational model to be well described like this:
	
	$$ C = n \cdot G^2 \cdot f \cdot t $$
	
	Where:
	
	\begin{itemize}
		\item $C$ is the total number of operations necessary to perform the simulation
		
		\item $n$ is the number of events that happen in or to the cell in one frame
		
		\item $G$ is the number of genes the cell carries. Squaring it just means that every interaction of two genes is considered, including the interaction of every gene with itself.
		
		\item $f$ is the frequency of updates, ie. "frames per second" of the simulation
		
		\item $t$ is the length of the simulation.
	\end{itemize}
	
	For some rough bounds on the numbers:
	
	\medskip
	
	The human genome contains around 19 000 to 20 000 protein-coding genes, around 13 000 inactive pseudogenes, and an number I couldn't determine of non-coding genes that don't produce protein yet are still critical to celular function. This adds up to a total of 32 000 to 33 000 genes plus the non-coding ones, I'll assume a total of around \textbf{45 000}. (\href{https://en.wikipedia.org/wiki/Human_genome#Molecular_organization_and_gene_content}{source: wikipedia})
	
	\medskip
	
	Transcription of genes happens surprisingly slowly in cells: "RNA polymerase builds the pre-mRNA molecule at a rate of 20 nucleotides per second enabling the production of thousands of pre-mRNA molecules from the same gene in an hour." - Based on this and my intuition, as a rough approximation I'll assume that a framerate of one frame per microsecond is sufficient. (\href{https://en.wikipedia.org/wiki/Protein_biosynthesis#Transcription}{source: wikipedia})
	
	\medskip
	
	Number of events involving the cell that might happen within one frame is a parameter I have no idea how to pin down, so I will arbitrarily assume it to be 1 as that gives a result which can later be simply multiplied if a new reasonable value is found.
	
	\medskip
	
	Time of simulation is easier to approximate but I would expect it to vary wildly per use case, so I will leave it at one second for now for similar reasons.
	
	$$ 1 \cdot (45 000)^2 \cdot (1\mu s)^-1 \cdot 1s = 2.025 \cdot 10^{15} $$
	
	Given these assumptions, simulation of one cell requires 2.025e15 operations per second. Given that there are about 40 trillion cells in the human body, a full-body simulation would require about 8.1e28 operations per second.
	
	\medskip
	
	\noindent\rule{\linewidth}{1pt}
	
	\medskip
	
	Roughly assuming 30 000 bit flips per operation to bring it into units more commonly used in the Stareater Expanse setting, that would translate to flip rates of 6.075e19 bits/s per cell, and 2.43e33 bits/s for a full body simulation.
	
	\medskip
	
	Memory use would be more complicated but I expect it somewhere between 720 kilobits per cell (16 bits per gene) and 12.4 gigabits (4 bits per base pair) per cell. For a full body that is between 28.8e19~bits and 4.96e23 bits.
	
	\medskip
	
	In either case as of the current state of computational hardware in the Stareater Expanse setting, for full body simulations those are computational requirements warranting a large computing cluster, even more massive than those used to run uploaded human minds. It is still well within the bounds of what a person might privately own for personal use, but also well past the point where one may want to simply rent the computing power instead of building out a privately owned cluster to that level.
	
	\medskip
	
	Typical household computational hardware would be better suited for very local simulations involving perhaps a million cells or less, but the computational capacity of cybernetic implants would not be sufficient for pretty much any useful level of cellular simulation without cooking the owner with waste heat in the process.
	
	\pagebreak
	
	\section{Molecular-level abstraction}
	
	"Molecular" would have to actually mean atomic, just without considering the internal states of atoms, since molecules aren't cleanly separable objects, they bind together, break apart, gain and lose tiny pieces occasionally, which would likely be even more computationally expensive to abstract away than just giving up and tracking every single atom and its interaction with the closest atoms around it.
	
	\medskip
	
	\noindent A basic approximation of the amount of computation needed for a given simulation would be as such:
	
	$$ C = a \cdot d \cdot i \cdot f \cdot t$$
	
	where:
	
	\begin{itemize}
		\item $C$ is the total number of operations
		
		\item $a$ is the number of atoms in the simulation
		
		\item $d$ is the number of atoms that any given atom is interacting with at one time (ie. number of neighbours it is in some type of contact with)
		
		\item $i$ is the number of different interactions that have to be considered every time two atoms are in contact
		
		\item $f$ is the framerate of the simulation
		
		\item $t$ is the length of the simulation
	\end{itemize}
	
	\textbf{bounds on $a$:} Given that the topic here is biology and cellular simulations, most of the human body could likely be fairly well approximated by water in terms of how many atoms there are in any given volume. 1 mole of water is about 18.015 milliliters and contains 6.022e23 molecules, which is 18.066e23 atoms. This means that there are roughly 1.003e23 atoms in each cubic centimeter of water.
	
	\medskip
	
	Here is a table with cell volumes sourced from [\href{https://book.bionumbers.org/how-big-is-a-human-cell/}{source}] with added atom numbers calculated using the approximation of 1.003e23 atoms per cm$^3$:
	
	\begin{center}
		\begin{tabular}{|l|c|c|}
			\hline
			\textbf{cell type} & \textbf{average volume (\textmu m$^3$)} & \textbf{num. of atoms} \\
			\hline
			sperm cell & 30 & 3.009e12 \\
			red blood cell & 100 & 1.003e13 \\
			lymphocyte & 130 & 1.3039e13 \\
			neutrophil & 300 & 3.009e13 \\
			beta cell & 1 000 & 1.003e14 \\
			enterocyte & 2 000 & 2.006e14 \\
			fibroblast & 3 000 & 3.009e14 \\
			HeLa, cervix & 4 000 & 4.012e14 \\
			hair cell (ear) & 4 000 & 4.012e14 \\
			osteoblast & 5,000 & 5.015e14 \\
			alveolar macrophage & 5 000 & 5.015e14 \\
			cardiomyocyte & 15 000 & 1.5045e15 \\
			megakaryocyte & 30 000 & 3.009e15 \\
			fat cell & 600 000 & 6.018e16 \\
			oocyte & 4 000 000 & 4.012e17 \\
			\hline
		\end{tabular}
	\end{center}
	
	For a whole-body simulation (even though it is not an expected use case) I would assume a 75L of volume, and correspondingly 7.5225e27 atoms.
	
	\pagebreak
	
	\textbf{bounds on $d$:} It may be expected for an atom to influence in some meaningful way every other atom with which it forms a molecule. The average mass of molecules in eukaryotic organisms is 49 kilodalton (\href{https://en.wikipedia.org/wiki/Protein#Biosynthesis}{source, "The average size of a protein increases from Archaea to Bacteria to Eukaryote (283, 311, 438 residues and 31, 34, 49 kDa respectively)"}).
	
	\medskip
	
	Without detailed analysis, 49kDa is roughly the mass of 4084 carbon atoms, so I will approximate $d$ as 4084.
	
	\medskip
	
	\textbf{bounds on $i$:} here are distinct influences I can come up with for each atom: atom's own innertia, electrostatic attraction to nearby atoms, electric attraction or repulsion between charged atoms. As such, without knowing any better, I'm approximating $i$ to be 3.
	
	\medskip
	
	\textbf{bounds on $f$:} Having ran the "folding at home" software to volunteer my computing power to medical research, I noticed the time length of a single simulation frame is given in femtoseconds, so I will set $f$ to 1fs$^-1$
	
	\medskip
	
	\textbf{bounds on $t$}: protein can take whole milliseconds to fold into the right shape, which is the low end of what could be usefully achieved with such a simulation. I fail to see any upper bound, it will likely vary wildly with use case, but the time factor can be simply removed in favour of calculating operations per second instead of total operations so that's what we'll do.
	
	\medskip
	
	Finally this results in an equation with one variable, the number of atoms or alternatively volume of biological tissue:
	
	$$ C = a \cdot 4084 \cdot 3 \cdot (1fs)^-1 = a \cdot 1.2252e19 Hz $$
	
	$$ a = V \cdot \frac{1.003e23}{1ml} = V \cdot 1.003e29m^{-3} $$
	
	where:
	
	\begin{itemize}
		\item $C$ is the computational power needed (operations per second) for a realtime simulation
		
		\item $a$ is the number of atoms in the simulation
		
		\item $V$ is the volume of tissue to be simulated
	\end{itemize}
	
	\medskip
	
	One operation will also be treated as 30 000 bit flips for the sake of bringing it closer to the units commonly used in the Stareater Expanse setting.
	
	\medskip
	
	Here are some relevant results:
	
	\begin{center}
		\begin{tabular}{|c|c|c|c|}
			\hline
			object of simulation & number of atoms & computing power & flip rate \\\hline
			per 1 atom on average & 1 & 1.2252e19 op/s & 3.6756e23 bits/s \\
			per 1ml of biological tissue & 1.003e23 & 1.2288756e42 op/s & 3.6866268e46 bit/s \\\hline
			ribosome & 300 000 & 3.6756e24 op/s & 1.10268e29 bits/s \\
			common cold virus (0.001$\mu m^3$) & 100 300 000 & 1.12288756e27 op/s & 3.6866268e31 bits/s \\
			mitochondrion (9$\mu m^3$) & 9.027e11 & 1.10598804e31 op/s & 3.31796412e35 bits/s \\
			fibroblast & 3.009e14 & 3.6866268e33 op/s & 1.10598804e38 bits/s \\
			oocyte & 4.012e17 & 4.9155024e36 op/s & 1.47465072e41 bits/s \\
			human embryo (week 8, 16ml) & 1.6048e24 & 1.96620096e43 op/s & 5.89860288e47 bits/s \\
			average adult human body (75L) & 7.5225e27 & 9.216567e46 op/s & 2.7649701e51 bits/s \\\hline
		\end{tabular}
	\end{center}
	
	For a bit of context, an analysis by Isaac Arthur [\href{https://www.youtube.com/embed/Ef-mxjYkllw?t=1116}{source}] puts the flip rate of a full matrioska brain built around the sun at around 1.45e49 bits/s, which is notably 2 orders of magnitude smaller than my estimation of what would be required to simulate a full human body.
	%Which may feel like a bit of relief for those who noticed how morally horrifying the last 2 options listed in the table are given the level of precision of the simulation likely being sufficient to simulate neurological activity.
	%Naturally some others will conclude that it's more practical to simply test outside of a simulation instead.
	
\end{document}